\documentclass[answers,12pt,addpoints]{exam} 
\usepackage{import}
\usepackage{graphicx}

\import{C:/Users/prana/OneDrive/Desktop/MathNotes}{style.tex}

% Header
\newcommand{\name}{Pranav Tikkawar}
\newcommand{\course}{STAT 669}
\newcommand{\assignment}{Assignment 1.2}
\author{\name}
\title{\course \ - \assignment}

\begin{document}
\maketitle

\begin{center}
    \includegraphics[width=0.8\textwidth]{"img/contour_plot.png"}
\end{center}
Above is the contour plot of a section of Death Valley, California. 
The specific data is from "USGS\_13\_n37w118\_20260112.tif". 
Some key features I am noticing are the extreme elevation changes in the valley area,
as well as the steep gradients in the surrounding mountainous regions.
The low elevation areas have large contour areas, while the high elevation areas have tightly packed contours. I was surprised to see how quickly the elevation changes in certain areas even for an area that is know for have its odd terrain.



\end{document}